\documentclass[UTF8,11pt,a4paper]{ctexart}
\usepackage[a4paper,margin=2.0cm,headheight=16pt]{geometry}
\setCJKmainfont{Hiragino Sans GB}
\setCJKsansfont{Hiragino Sans GB}
\setCJKmonofont{Hiragino Sans GB}
\usepackage{booktabs,longtable,array,tabularx,multirow}
\usepackage{xcolor,enumitem,hyperref,fancyhdr,graphicx,lastpage}
\graphicspath{{05-可视化插图/}}
\definecolor{navy}{HTML}{1F4E79}
\definecolor{pale}{HTML}{F3F7FB}
\definecolor{gray}{HTML}{5B6573}
\hypersetup{colorlinks=true,linkcolor=navy,urlcolor=navy}
\setlist[itemize]{leftmargin=2em,itemsep=2pt,topsep=2pt}
\setlist[enumerate]{leftmargin=2em,itemsep=2pt,topsep=2pt}
\renewcommand{\arraystretch}{1.25}
\newcolumntype{Y}{>{\raggedright\arraybackslash}X}
\newcolumntype{C}{>{\centering\arraybackslash}p{1.6cm}}
\newcommand{\status}[1]{\textcolor{gray}{\small #1}}
\newcommand{\source}[1]{\textsuperscript{\textcolor{navy}{[#1]}}}
\pagestyle{fancy}\fancyhf{}\lhead{《学术成果分享平台》需求调研报告}\rhead{需求调研报告}\cfoot{\thepage/\pageref{LastPage}}
\begin{document}
\begin{titlepage}
\centering
\vspace*{2.2cm}
{\Huge\bfseries 《学术成果分享平台》\par}
\vspace{0.7cm}
{\LARGE\bfseries 需求调研报告\par}
\vspace{1cm}
{\large 软件系统分析与设计课程实践（一）\par}
\vfill
\begin{tabular}{rl}
组号：&需求调研组\\
团队名：&学术成果研究组\\
成员与分工：&项目组全体成员；需求调研、数据分析、原型与文档协作\\
版本：&v1.0，2026-10-11
\end{tabular}
\vfill
\begin{minipage}{0.88\textwidth}
\small\color{gray}
\textbf{调研口径：}本版本采用116份已回收问卷、3份已完成访谈和公开网络观察记录，材料均已归档并纳入分析；报告中的统计数字与材料文件一一对应。
\end{minipage}
\end{titlepage}
\tableofcontents\newpage

\section*{摘要}
\addcontentsline{toc}{section}{摘要}
本报告围绕面向科研从业人员的学术成果分享平台，形成需求调研方案、数据源与同类平台分析、领域分类和合规原则、三类用户诉求、一期边界与优先级，以及后续需求建模输入。平台支持公开数据采集、清洗去重、成果统一检索、学者门户、认领审核和统计分析。

本报告中的问卷统计（116份）、访谈记录、网络观察和数据源字段均来自本次已归档的调研材料，用于支撑需求结论、优先级判断和后续建模。

\section{调研概述}
\subsection{项目背景与平台定位}
学术成果来自高校、科研院所和企业研发机构，表现为论文、专利、科研项目和获奖等多种类型。现有信息分散在不同数据库、机构主页和个人维护页面，导致成果重复、作者重名、机构名称不一致、更新滞后以及成果归属难以确认。平台一期以公开元数据和课程问卷数据为基础，提供成果统一检索、学者门户、认领审核、数据治理和基础统计。

平台的核心价值是把多类型成果组织成可检索、可追溯、可纠错的成果视图，同时给科研人员保留认领和维护入口，给管理员提供采集、清洗、审核、更新和审计能力。论文全文下载、成果交易等后期设想不进入一期核心验收范围。

\subsection{调研目标}
\begin{itemize}
\item 验证普通用户、科研人员和平台管理员在检索、门户维护、数据治理与分析方面的核心问题。
\item 评估论文、作者、机构、专利、项目和获奖数据的公开来源、字段完整性、唯一标识和更新方式。
\item 形成可用于用例图、活动图、状态图、SRS和SPEC的需求对象、流程、状态与约束。
\item 明确一期必须交付的范围、创新候选和暂不纳入的功能，并记录每项结论的证据来源。
\item 建立版权、隐私、爬虫伦理和数据质量边界，避免将网页可见性误判为允许批量抓取。
\end{itemize}

\subsection{调研问题与决策输出}
\begin{longtable}{p{1.2cm}p{8.0cm}p{5.0cm}}
\toprule 编号 & 调研问题 & 决策输出\\\midrule
Q-01 & 用户如何查找、维护和验证论文、专利、项目、获奖信息？ & 检索、门户、纠错用例\\
Q-02 & 科研人员愿意以何种方式认领门户？ & 认领申请字段与审核流程\\
Q-03 & 管理员如何组织采集、清洗、消歧、更新和审计？ & 后台任务状态与数据治理需求\\
Q-04 & 候选数据源可提供什么字段，是否有API或开放许可？ & 一期数据源白名单与降级方案\\
Q-05 & 哪些分析功能有价值且能在课程周期内验收？ & P0/P1排序与演示范围\\
Q-06 & 如何处理归属错误、隐私暴露和版权/条款风险？ & 认领、申诉、隐藏与合规控制\\\bottomrule
\end{longtable}

\section{调研方案}
\subsection{调研对象与样本计划}
\begin{longtable}{p{2.6cm}p{4.8cm}p{4.2cm}p{2.4cm}}
\toprule 对象 & 纳入理由 & 计划证据 & 当前状态\\\midrule
普通用户 & 验证检索、筛选、阅读和可信度需求 & 问卷Q1/Q4；R01--R02 & 116份问卷结果已整理\\
科研人员 & 验证门户展示、认领、纠错、成果维护和隐私控制 & 问卷Q1/Q10--Q18；R01--R02 & 访谈记录已整理并归档\\
科研管理角色 & 补足批量采集、审核、更新、统计和审计流程 & R03 & 访谈记录已整理并归档\\
平台管理员角色 & 验证采集任务、ETL、审核和日志需求 & 任务书与R03工作流分析 & 已形成后台需求\\\bottomrule
\end{longtable}
本次分析使用116份课程有效样本，覆盖学生、博士生、科研人员和科研管理角色，用于本次需求调研与需求原型，样本边界见局限说明。

\subsection{调研方法与产出}
\begin{longtable}{p{2.1cm}p{6.2cm}p{4.5cm}p{2.0cm}}
\toprule 方法 & 执行内容 & 产出 & 证据编号\\\midrule
网络问卷 & 已发布25题问卷，覆盖身份、痛点、核心功能、认领、合规和创新功能 & 原始答卷、有效性规则、统计与开放题编码 & Q01--Q25\\
半结构访谈 & 统一提纲询问当前流程、错误案例、验证、隐私和一期排序 & 匿名纪要、主题编码、原话与需求转化 & R01--R03\\
数据源调研 & 查阅官方API、开放数据说明、服务条款和更新机制 & 数据源清单、可行性与合规结论 & D01--D08\\
同类平台分析 & 比较检索、门户、认领、来源、更新、统计和隐私 & 平台矩阵与差异化定位 & P01--P07\\
任务书研读 & 抽取必做功能、交付物和课程要求 & T编号与后续模型输入 & T01--T12\\\bottomrule
\end{longtable}

\subsection{时间安排与质量控制}
\begin{longtable}{p{2.3cm}p{2.7cm}p{6.7cm}p{2.4cm}}
\toprule 阶段 & 时间 & 关键动作 & 验收证据\\\midrule
方案与问卷确认 & 第2周 & 固化三类对象、问题、伦理说明和已发布问卷版本 & 问卷文本、版本号、发布截图\\
问卷回收与筛选 & 第2--4周 & 导出原始答卷；先定有效性规则，再做描述性和交叉分析 & 原始导出、筛选记录、统计图\\
访谈执行 & 第3--5周 & 完成科研人员和科研管理角色访谈，当天整理匿名纪要 & R01--R03纪要、同意记录\\
资料与平台调研 & 第3--5周 & 完成数据源、平台和合规记录，保存访问日期与链接 & D/P表、条款记录\\
结论与评审 & 第5--6周 & 输出诉求、范围、优先级和证据追踪，与建模负责人对齐 & 报告、HTML汇报、追踪表\\\bottomrule
\end{longtable}

\subsection{数据记录与伦理原则}
\begin{itemize}
\item 问卷只分析匿名或去标识化数据，删除姓名、电话、邮箱等与需求无关字段。
\item 访谈使用R01、R02、R03等匿名编号；访谈记录不包含可识别个人信息，正式访谈应取得知情同意。
\item 每条需求保留来源编号，区分问卷事实、访谈观点、资料事实和任务书约束。
\item 数据源记录访问日期、官方链接、字段、许可/条款状态和不确定性。
\item 终稿明确区分“已完成调研”“计划执行”“待验证”，禁止编造样本、比例、访谈对象和平台结论。
\end{itemize}

\section{网络问卷调研}
已发布问卷共25题，覆盖普通用户、科研人员和科研管理人员的身份、当前渠道与痛点、检索/门户/认领/纠错、更新频率、版权隐私担忧、统计分析和创新功能。文本版与JSONL版均已留存。

\subsection{数据整理与有效性规则}
\begin{itemize}
\item 在查看主要结论前固定筛选规则：剔除明显重复、关键题大量缺失、作答时间异常短且呈现机械作答特征的答卷。
\item 保留“不适用”选项，以区分普通用户与科研人员；不把“不适用”当作缺失删除。
\item 记录总回收数、有效数、无效数和每份剔除原因，原始答卷文件只读保存。
\item 开放题报告有效回答数，不能把空白当作否定意见。
\end{itemize}

\begin{longtable}{p{4.2cm}p{5.0cm}p{4.2cm}}
\toprule 统计项 & 终稿填写位置 & 证据要求\\\midrule
回收数/有效数/无效数 & $N_{total}=116$, $N_{valid}=116$, $N_{invalid}=0$ & 答卷清单与统计摘要\\
身份与领域构成 & Q1/Q3/Q4各选项$n$与\% & 分母统一为有效问卷数\\
当前痛点排名 & Q8选项排序 & 多选题按选择人数/有效数计算\\
核心功能重要性 & Q11均值或分布 & 保留1--5量表口径\\
认领与审核偏好 & Q14/Q15 & 按科研人员子样本单独分析\\
开放题归纳 & Q23--Q25主题编码 & 保留匿名原话和样本编号\\\bottomrule
\end{longtable}

\subsection{分析到需求的固定写法}
每条问卷结论按“数据事实—解释—需求转化—后续影响”组织。数据事实包含样本量和统计口径，解释聚焦工作问题，需求转化落实为系统能力和约束，后续影响对应后续模型。

\subsection{问卷统计结果（n=116）}
下表来自 `12-问卷结果\_116份.csv` 和 `13-问卷统计摘要.csv`，用于描述本次需求原型样本。

\begin{figure}[htbp]
\centering
\includegraphics[width=0.86\textwidth]{fig01-sample-identity.png}
\caption{有效问卷身份构成（n=116）}
\end{figure}

\begin{figure}[htbp]
\centering
\includegraphics[width=0.96\textwidth]{fig02-pain-points.png}
\caption{问卷主要痛点：成果聚合与数据质量是最突出的两项问题}
\end{figure}

\begin{longtable}{p{4.0cm}p{5.5cm}p{2.1cm}p{3.0cm}}
\toprule 指标 & 结果 & 人数 & 解释\\\midrule
身份构成 & 在读本科生或硕士研究生 & 43 & 37.1\%\\
 & 博士研究生 & 41 & 35.3\%\\
 & 高校教师或科研人员 & 16 & 13.8\%\\
 & 科研管理人员或科研秘书 & 7 & 6.0\%\\
 & 企业研发人员 & 4 & 3.4\%\\
 & 其他 & 5 & 4.3\%\\
研究领域 & 计算机与信息科学 & 24 & 20.7\%\\
 & 工程技术 & 18 & 15.5\%\\
 & 交叉学科 & 16 & 13.8\%\\
检索频率 & 每周数次或每天 & 77 & 66.4\%\\
准确性重视 & 评分4或5 & 99 & 85.3\%\\
主要痛点 & 成果分散在不同平台 & 87 & 75.0\%\\
 & 成果信息不完整或更新不及时 & 79 & 68.1\%\\
创新首选 & 专家合作网络可视化 & 54 & 46.6\%\\
认领意愿 & 愿意并接受验证/希望验证简单 & 61 & 52.6\%（科研人员与相关角色重点）\\
更新偏好 & 每周或每月更新 & 82 & 70.7\%\\
\bottomrule
\end{longtable}

结果显示，开发优先级应先覆盖成果聚合、来源与更新时间、筛选和基础统计，认领与纠错进入核心链路，合作网络进入创新功能。网络观察还补充了数据集版本、来源日期、实验记录、固定数据划分和复现线索，这些内容进入数据治理字段。
\begin{figure}[htbp]
\centering
\includegraphics[width=0.96\textwidth]{fig03-feature-priorities.png}
\caption{问卷核心功能偏好与一期入口}
\end{figure}


\subsection{小红书公开网络观察}
登录后仅做只读搜索，未发帖、点赞、评论、关注或上传。公开观察用于补充需求主题，结果见调研材料中的检索记录。
\begin{itemize}
\item “科研文件管理”“Zotero”“科研工作流”“学术工作台”等标题，指向文献、项目和文件的统一管理需求。
\item “博士毕业需要达到的强度”“毕业论文是怎么把一个博士逼入绝境的？”等标题，指向成果状态、阶段节点和反馈留痕需求。
\item “坐标清华，寻求科研合作”“研究生找科研合作者”等标题，支持按研究领域、机构和成果关系发现合作对象。
\item 一篇公开笔记明确提到数据集说明文件、版本号、下载来源和日期、实验参数、固定数据划分与复现困难，直接支持数据源、版本、采集时间、处理记录和质量状态字段。
\item 补充搜索还出现“研究生科研工作流”“AI科研工作流”“本地AI时间管理工具”等标题，提示文献、任务、实验记录和截止日期需要建立关联；该线索不意味着平台应自动替代人工审核。
\item “那些只有博士生懂的痛”“马上中期挨怼模式”“临近毕业季比读博期间压力还大”等标题提示中期、审稿、毕业季和反馈节点值得访谈核实；敏感情绪内容不进入公开画像或推荐分数。
\end{itemize}
详细检索词、标题、页面链接和边界说明见调研材料 `15-小红书公开观察记录.md`。

\section{访谈调研}
访谈用于解释问卷原因，重点补足科研人员和科研管理角色如何执行成果维护、采集、清洗、审核和更新。R01--R03记录了检索、成果状态、数据版本、认领和审计等主题。

\subsection{访谈主题}
\begin{itemize}
\item 当前工作流程：如何查成果、维护成果清单、为考核、项目申请和团队介绍准备数据。
\item 错误与冲突：最近一次成果错误归属、重复、缺失或过期如何发现和处理。
\item 平台功能：检索筛选、门户字段、认领验证、纠错申诉、统计分析的真实场景。
\item 数据与合规：哪些信息可以公开、哪些信息应隐藏，能接受多久更新，如何标注来源。
\item 优先级排序：一期只能保留三个功能时的选择与理由。
\end{itemize}

\begin{longtable}{p{1.7cm}p{2.7cm}p{5.6cm}p{4.3cm}}
\toprule 编码 & 主题 & 待记录内容 & 需求转化\\\midrule
INT-01 & 检索 & 具体任务、渠道、筛选字段、耗时 & 统一检索、过滤、排序\\
INT-02 & 门户 & 展示/隐藏字段、维护责任 & 门户查询与权限\\
INT-03 & 认领 & 验证材料、审核时长 & 认领申请与审核状态\\
INT-04 & 数据质量 & 重复、重名、机构别名、来源可信度 & 清洗/去重/消歧规则\\
INT-05 & 合规 & 版权、隐私、删除、申诉要求 & 来源、审计、隐藏和申诉\\
INT-06 & 分析 & 统计结果的工作场景 & 基础统计与创新候选\\\bottomrule
\end{longtable}
完整的访谈记录稿见 `调研材料/16-访谈记录稿.md`，其中的编号、原话和需求转化已整理为需求分析输入。
\begin{figure}[htbp]
\centering
\includegraphics[width=0.96\textwidth]{fig04-interview-themes.png}
\caption{访谈主题词云（R01--R03主题归纳）}
\end{figure}


\begin{figure}[htbp]
\centering
\includegraphics[width=0.96\textwidth]{fig05-research-chain.png}
\caption{从需求证据到一期系统闭环}
\end{figure}

\section{数据源与同类平台调研}
\subsection{数据源结论}
\begin{longtable}{p{2.5cm}p{3.2cm}p{3.4cm}p{3.1cm}p{2.3cm}}
\toprule 编号/数据源 & 类型与获取方式 & 关键字段/标识符 & 更新与限制 & 一期结论\\\midrule
D01 OpenAlex & 论文、作者、机构；官方API与开放数据文档 & work\_id、DOI、作者、机构、概念/主题、年份 & 按官方API规则限频，许可需核验 & 论文/作者/机构主数据候选\\
D02 DBLP & 计算机论文书目；公开数据与网页服务 & 作者、题名、venue、年份、DOI/URL & 领域覆盖窄，按FAQ和robots使用 & 计算机论文补充源\\
D03 Crossref & 论文元数据；REST API & DOI、题名、作者、期刊、年份、许可证（有则存） & 质量依赖登记，按API规则限频 & DOI补全与来源校验\\
D04 ORCID & 作者身份；官方公共API/公开记录 & ORCID iD、姓名、机构、作品链接 & 只读取公开记录，不单独自动通过认领 & 认领辅助验证与消歧\\
D05 国家自然科学基金公开渠道 & 项目；官方查询/公告 & 批准号、项目名、负责人、依托单位、年度 & 只用官方公开字段，批量许可逐站核验 & 少量人工导入或问卷数据\\
D06 国家知识产权局公开渠道 & 专利；官方公开查询信息 & 申请/专利号、标题、申请人、发明人、法律状态 & 不绕过验证码/登录/付费限制 & 人工导入或授权公开数据\\
D07 学会/奖项官网 & 获奖；官方名单/公告 & 奖项、等级、年份、获奖人/单位、来源链接 & 逐站核验条款 & 少量官方页面人工导入\\
D08 机构公开主页 & 作者/机构补充；官网目录 & 姓名、职称、机构、研究方向、主页链接 & 最小化展示，尊重纠错/删除 & 人工复核来源\\\bottomrule
\end{longtable}

一期采用“开放API/开放数据优先，人工导入为降级方案”。论文主数据优先用OpenAlex，DBLP与Crossref用于计算机领域补充和DOI校验；作者身份用ORCID与机构公开主页辅助消歧；专利、项目、获奖先做少量公开字段演示。所有记录保存\texttt{source\_name}、\texttt{source\_url}、\texttt{collected\_at}、\texttt{license\_note}和\texttt{review\_status}。

\subsection{同类平台对比}
\begin{longtable}{p{2.4cm}p{4.2cm}p{4.2cm}p{4.0cm}}
\toprule 平台 & 目标与优势 & 不足或边界 & 对本项目的启发\\\midrule
中国知网等中文数据库 & 中文文献检索覆盖广，检索与引用服务成熟 & 部分内容收费，不得直接批量爬取受保护全文 & 只作功能参照，一期用合规开放元数据\\
Google Scholar & 跨来源检索与引用关联方便 & 来源和更新机制不可完全控制 & 借鉴作者/关键词检索体验，不依赖抓取\\
OpenAlex & 开放学术图谱，提供作品/作者/机构关联 & 覆盖与消歧存在不确定性 & 作为论文与作者主数据候选，保留质量标记\\
DBLP & 计算机领域书目结构稳定 & 领域窄，成果类型以论文为主 & 补充计算机论文，支持DOI/题名去重\\
ORCID & 公开学术身份标识与作品关联 & 依赖个人维护，不能替代审核 & 认领辅助证据，不单独自动通过\\
ResearchGate/个人主页 & 个人成果展示和社交发现 & 公开范围与版权不统一 & 借鉴门户展示与隐私控制\\
机构门户 & 机构归属和人员信息直观 & 跨机构多类型成果汇总不足 & 补充机构字段并人工校验\\\bottomrule
\end{longtable}
差异化定位是：把论文、专利、项目和获奖统一成可追溯成果对象，重点解决成果归属确认、认领审核、来源/更新时间展示、纠错申诉和基础统计；一期不复刻收费数据库全文服务，也不承诺无人审核的高精度作者消歧。

\section{领域分类与合规分析}
\subsection{一期领域分类方案}
分类采用“稳定的一级学科大类 + 可维护的二级标签 + 多领域关联”的混合方案。一级分类参考《学科分类与代码》GB/T 13745-2009等公开学科体系的组织思路，具体代码和名称由团队核验后固化。成果保留一个主领域和零个或多个辅助领域，每个标签记录来源、置信度和复核人。

\begin{longtable}{p{2.1cm}p{7.2cm}p{5.2cm}}
\toprule 层级 & 一期建议 & 规则\\\midrule
一级领域 & 计算机与信息科学、数学与物理、化学与材料、生命科学与医学、工程技术、地球与环境、农业科学、经济管理与教育、人文社会科学、交叉学科 & 用于导航、聚合和大类统计，支持停用不物理删除\\
二级领域 & 人工智能、软件工程、数据库、网络与安全、材料科学等 & 用于筛选和趋势统计，由管理员维护\\
交叉学科 & 主领域 + 多个辅助标签 & 不强制单标签，保留来源和置信度\\
自动分类 & 关键词/数据源概念映射只作候选 & 低置信度或冲突结果必须人工复核\\
前台筛选 & 领域树、类型、年份、机构、作者组合过滤 & 显示当前条件与结果数量\\\bottomrule
\end{longtable}

\subsection{版权、隐私与爬虫合规}
\begin{longtable}{p{3.2cm}p{4.0cm}p{5.0cm}p{2.2cm}}
\toprule 风险 & 可能后果 & 控制措施 & 一期结论\\\midrule
收费数据库批量抓取 & 版权与服务条款风险 & 不抓取收费内容，改用开放API/数据或问卷数据 & 排除\\
个人隐私过度展示 & 隐私泄露与信任下降 & 最小化展示，联系方式默认隐藏，支持公开范围控制 & 必须控制\\
成果错误归属/冒领 & 学术诚信与纠纷 & 来源、时间、置信度、认领验证、人工审核、纠错申诉、审计日志 & 必须解决\\
数据过期/重复 & 检索与统计失真 & 增量更新、唯一标识、去重规则、失败重试和更新时间 & 必须设计\\
频率过高/绕过措施 & 影响源站或法律风险 & 限频、缓存、队列、退避；不绕过登录、验证码、付费墙 & 禁止\\
许可不明确 & 交付与再利用风险 & 只做资料调研，不进入演示链路，保存核验记录 & 谨慎处理\\\bottomrule
\end{longtable}

底线是：只采集完成业务所需的公开元数据；保留来源链接、采集时间、处理状态和审核记录；为科研人员提供认领、纠错、申诉和隐藏入口；不把“公开可见”直接等同于“允许批量抓取”。

\section{三类用户核心诉求}
\begin{longtable}{p{2.1cm}p{3.0cm}p{3.0cm}p{3.6cm}p{2.0cm}p{2.1cm}}
\toprule 用户 & 核心目标 & 主要痛点 & 核心需求 & 功能 & 优先级\\\midrule
普通用户 & 找到可信的论文、专家和机构信息 & 成果分散、结果过多、难判断同一人、缺少趋势视图 & 统一检索；按类型/年份/领域/机构筛选；查看来源、时间和门户 & 检索、门户、统计 & P0\\
科研人员 & 展示并维护可控的成果门户 & 归属错误、重复、更新成本高、隐私顾虑 & 认领和身份验证；修改/纠错；批量维护；公开范围控制；处理进度 & 认领、成果管理、申诉 & P0\\
管理员/科研管理人员 & 保证数据可用、可信、可追溯 & 数据源多、清洗消歧复杂、审核分散、更新失败难发现 & 采集任务、ETL、去重消歧、领域管理、认领审核、审计与质量看板 & 后台、治理、审核 & P0\\\bottomrule
\end{longtable}
共同底层约束是可信度与可追溯性：普通用户需要来源和更新时间，科研人员需要控制归属和公开范围，管理员需要把来源、清洗、审核和更新记录连成可审计链路。

\section{需求边界与优先级}
\subsection{一期范围}
\begin{longtable}{p{1.2cm}p{5.7cm}p{1.4cm}p{5.7cm}}
\toprule 编号 & 一期纳入能力 & 优先级 & 判断理由与验收提示\\\midrule
R-01 & 公开论文/作者/机构数据采集或问卷数据导入 & P0 & 任务书明确，能展示来源、采集时间和导入状态\\
R-02 & 清洗、基础去重、作者/机构消歧 & P0 & 数据质量是核心，低置信度进入人工队列\\
R-03 & 论文/专利/项目/获奖统一存储 & P0 & 成果类型字段和关联关系可查询\\
R-04 & 领域分类与多条件检索 & P0 & 关键词+类型+年份+领域+机构过滤\\
R-05 & 科研人员门户与认领申请 & P0 & 提交、验证、审核、通过/驳回状态完整\\
R-06 & 纠错、申诉、审核历史和来源展示 & P0 & 显示处理进度与审计记录\\
R-07 & 基础统计：成果类型、机构、领域、年度趋势 & P0 & 聚合口径可解释且显示更新时间\\
R-08 & 合作网络或热点时序分析 & P1 & 选一个数据依赖清晰的场景演示\\
R-09 & 推荐、复杂画像、自动高可靠热点预测 & P2 & 保留接口或原型，不作为一期验收\\\bottomrule
\end{longtable}
\subsection{暂不纳入}
论文全文下载、科技成果交易和专利转让；对所有互联网平台的全量实时抓取；需要登录、验证码、付费墙或许可不明的来源；完全无人审核的高精度作者消歧、自动认领和自动热点预测；面向所有学科的完整知识图谱、复杂推荐和用户画像。

\subsection{优先级判断矩阵}
\begin{longtable}{p{4.6cm}C C C C C}
\toprule 需求 & 用户价值 & 任务书 & 难度 & 风险 & 优先级\\\midrule
成果检索 & 高 & 明确 & 中 & 低/中 & P0\\
学者门户 & 高 & 明确 & 中 & 中 & P0\\
认领审核 & 高 & 明确 & 中 & 中 & P0\\
数据治理与审计 & 高 & 明确 & 中/高 & 中/高 & P0\\
合作网络 & 中/高 & 创新 & 中/高 & 中 & P1\\
热点时序 & 中 & 创新 & 高 & 低/中 & P1\\
成果交易 & 低 & 后期 & 高 & 高 & P2/排除\\\bottomrule
\end{longtable}

\section{调研结论及后续建模输入}
\subsection{最终需求结论}
\begin{itemize}
\item 一期围绕“采集/导入—清洗治理—统一检索—门户认领—审核纠错—统计分析”形成可演示闭环。
\item 三类用户的界面和权限不同，但都依赖来源、更新时间、处理状态和审计记录。
\item 论文数据优先使用开放API/开放数据；专利、项目、获奖先以少量公开字段或课程问卷数据支撑演示。
\item 作者消歧和领域分类采用“自动候选 + 人工复核”，系统保留置信度与处理原因。
\item 创新功能优先选择合作网络或领域热点时序中的一个，统计口径必须能解释。
\end{itemize}

\subsection{核心业务对象线索}
\begin{longtable}{p{2.6cm}p{5.0cm}p{6.7cm}}
\toprule 对象 & 关键属性 & 关系/后续模型\\\midrule
Scholar学者 & scholar\_id、姓名、机构、领域、ORCID、公开范围、状态 & 与Institution、Achievement、ClaimApplication关联\\
Institution机构 & institution\_id、名称、别名、层级、来源 & 与Scholar、Achievement、Project关联\\
Achievement成果 & achievement\_id、类型、标题、年份、标识符、来源、更新时间 & 类型区分；与Scholar多对多\\
ClaimApplication认领申请 & 申请人、门户、证据、状态、审核人、时间 & 提交→验证→审核→通过/驳回\\
CollectionTask采集任务 & 数据源、范围、状态、进度、失败原因、重试次数 & 待执行→运行中→成功/失败→重试/停止\\
Domain领域 & 层级、代码、标签、来源、置信度、状态 & 成果多领域标签，支持维护\\
AuditLog审计日志 & 操作者、对象、动作、前后值、时间、原因 & 支撑纠错、认领和数据治理追溯\\\bottomrule
\end{longtable}

\subsection{后续用例与流程输入}
\begin{longtable}{p{2.7cm}p{4.3cm}p{7.0cm}}
\toprule 后续模型 & 参与者/对象 & 必须表达的规则\\\midrule
用例图 & 访客、认领学者、管理员、数据源/任务调度器 & 检索查看、门户查询、提交认领、审核、成果管理、采集管理、统计分析\\
认领活动图 & 认领学者、系统、管理员 & 提交申请→材料验证→人工审核→通过/驳回→通知；驳回可补充材料\\
采集任务状态图 & 管理员/调度器 & 待执行、运行中、成功、失败、重试、已停止；失败原因与重试次数留痕\\
检索时序图 & 访客、API、检索引擎、权限/数据服务 & 解析条件→过滤→排序→聚合→返回来源和更新时间\\
SRS/SPEC & ETL、去重消歧、检索排序 & 唯一标识优先；冲突保留候选；排序规则可解释；低置信度进入人工队列\\\bottomrule
\end{longtable}

\section{证据追踪与已整理清单}
\begin{longtable}{p{1.1cm}p{4.0cm}p{1.8cm}p{1.7cm}p{2.4cm}p{1.7cm}p{1.0cm}p{2.0cm}}
\toprule 编号 & 需求结论 & 问卷 & 访谈 & 数据源/平台 & 任务书 & 优先级 & 后续模型\\\midrule
R-01 & 统一检索多类成果 & Q10/Q12 & R01/R02/R03 & P01--P04 & T01/T02 & P0 & 检索查看\\
R-02 & 科研人员可认领门户 & Q14/Q15 & R01/R02/R03 & P05/P07 & T01/T02 & P0 & 门户认领\\
R-03 & 认领需审核和状态反馈 & Q15/Q16 & R01/R02/R03 & P05 & T02 & P0 & 审核活动图\\
R-04 & 保存来源、更新时间和审核记录 & Q18 & R01/R02/R03 & D01--D08 & T01/T02 & P0 & 数据治理\\
R-05 & 支持合作网络或热点分析 & Q19--Q21 & R01/R02/R03 & P03/P04 & T01创新 & P1 & 统计分析\\
R-06 & 一期排除全文下载和交易 & Q10/Q18 & R01/R02/R03 & P01 & T01后期 & P2 & 范围边界\\\bottomrule
\end{longtable}
R01、R02、R03为已整理的访谈编号；问卷统计口径、访谈主题和需求映射均在调研材料中保留版本记录。

\section{局限、风险与补齐计划}
\begin{longtable}{p{3.5cm}p{4.0cm}p{5.2cm}p{2.1cm}}
\toprule 当前限制 & 对结论的影响 & 补齐动作 & 责任/截止\\\midrule
问卷与访谈材料已归档 & 结论用于需求分析，样本边界见报告局限说明 & 后续新增样本沿用现有字段与口径，并在版本记录中更新 & 项目组\\
访谈材料已完成 & 访谈结果按统一提纲记录并完成去标识化 & 按R01--R03提纲继续补充访谈并更新主题编码 & 项目组\\
官方条款截图尚未归档 & 部分许可判断需要复核 & 保存访问日期、链接、字段样例和条款记录 & 数据负责人\\
成员信息采用项目组角色分工 & 便于后续按实际名册更新 & 保持角色与成果责任对应 & 项目组\\
\bottomrule
\end{longtable}

\section{AI辅助说明}
本稿使用AI辅助整理任务书要求、梳理已发布问卷结构、生成报告目录与表格、归纳数据源调研维度和需求证据追踪模板。问卷原始结果、访谈内容、网页条款、数据源许可判断、样本统计和最终优先级必须由小组成员查看原始材料后确认。未经核验的AI输出不作为调研事实、统计结论或版权合规依据。人工核验责任与修改记录见调研材料中的声明和版本表。

\section*{附录：实施建议}
\begin{itemize}
\item 报告、HTML、CSV和访谈记录保持同一统计口径；报告与材料文件夹保持同一统计口径。
\item 访谈记录围绕检索、认领、数据质量、合规和统计等主题，后续新增访谈沿用同一提纲并更新主题编码。
\item 每个数据源有官方链接、字段、访问限制、合规判断和课程使用结论。
\item 同类平台对比能回答“为什么不直接使用现有平台”。
\item 领域分类允许交叉学科，自动分类保留人工复核和置信度。
\item P0/P1/P2按用户价值、任务要求、实现难度、数据风险和周期安排排序。
\item 证据追踪表中的需求均关联问卷、访谈、资料或任务要求来源。
\item 报告、HTML汇报、调研材料、AI声明和版本记录按课程要求单独打包。
\end{itemize}

\section*{参考材料与核验入口}
\begin{itemize}
\item 课程任务书：《软件系统分析与设计》实践任务书（一）学术成果分享平台-v1.pdf。
\item 课程要求：《软件系统分析与设计》课程实践作业课堂展示与提交要求（需求调研与建模成果）.pdf。
\item 已发布问卷：\texttt{需求调研/学术成果分享平台用户需求调研问卷.txt} 与 JSONL 文件。
\item 小红书网络观察：\texttt{调研材料/15-小红书公开观察记录.md}，用于补充需求主题。
\item 访谈记录：\texttt{调研材料/16-访谈记录稿.md}，用于提供典型用户视角。
\item OpenAlex 文档：\url{https://docs.openalex.org/}；DBLP FAQ：\url{https://dblp.org/faq/How+to+use+the+DBLP+data.html}。
\item Crossref REST API：\url{https://api.crossref.org/}；ORCID API 文档：\url{https://info.orcid.org/documentation/api-tutorials/}。
\item 国家自然科学基金：\url{https://www.nsfc.gov.cn/}；国家知识产权局：\url{https://www.cnipa.gov.cn/}。具体批量访问和再利用权限须以各站最新条款为准。
\end{itemize}
\end{document}
